\begin{statement}{C.2.NC.1}{Known, with derivation}{The local-to-global Euler formula with multiplicities and self-intersections}
Let $g:U\to\Delta$ be a proper holomorphic map from a smooth complex threefold to a disc, smooth over $\Delta^*$, whose central fibre is $\sum_i m_iS_i$ and whose reduced support $S$ has normal crossings. Let $\nu:\coprod_i T_i\to S$ be normalization and let $D_i\subset T_i$ be its reduced conductor divisor. For a sufficiently small nonzero $t$,
\[
 e(U_t)=\sum_i m_i e(T_i\setminus D_i).
 \tag{C.NC.1}
\]
The formula does not require a Kähler total space, a polarization, uniform multiplicities, or simple normal crossings.
\end{statement}
\begin{proof}
We spell out the normal-crossings form of the construction in \cite[\S2.1--2.2, pp.~237--241]{ACampo1975}. Replace each local complex normal coordinate by its modulus and argument. In a chart with $r$ local branches, absorb a local root of the unit in the equation of $g$ into one coordinate, so that
\[
 g=z_1^{a_1}\cdots z_r^{a_r},\qquad
 z_j=\rho_j e^{\mathrm i\theta_j},\qquad 1\le r\le3.
\]
The resulting real oriented blowup $Z\to U$ has a boundary with corners over $S$. This construction glues even when branches are permuted: a change of defining equation multiplies a normal coordinate by a nowhere-zero holomorphic function, which multiplies its modulus by a positive function and translates its argument; permutation of the branches permutes the corresponding factors. Thus the construction uses the local branches themselves and does not require a global ordering or a global separation of the branches.

The lift of $g$ to the real oriented blowup of the disc is given in these coordinates by
\[
 (\rho_j,\theta_j)\longmapsto
 \left(\prod_j\rho_j^{a_j},\ \sum_j a_j\theta_j\pmod {2\pi}\right).
 \tag{C.NC.2}
\]
This map is locally a topological projection, including along the corners. For the radial variables this can be checked explicitly: put $u_j=\rho_j-\rho_r$ for $j<r$ and $v=\prod_j\rho_j^{a_j}$. Given $(u_1,\ldots,u_{r-1},v)$, the missing variable $\lambda=\rho_r$ is the unique solution, for $v>0$, of
\[
 \lambda^{a_r}\prod_{j<r}(u_j+\lambda)^{a_j}=v,
 \quad \lambda\ge\max(0,-u_1,\ldots,-u_{r-1});
\]
for $v=0$ it is the endpoint of that interval. Strict monotonicity for positive $v$ and continuity at the endpoint give the inverse radial chart. The angular coordinate in (C.NC.2) is a submersion in local real angular charts. Properness and the resulting local product charts give the boundary fibration used in the cited construction; its fibre over an argument $\theta$ is homeomorphic to $U_t$ for $\arg t=\theta$. This is exactly the comparison proved in the lemma of \cite[\S2.2]{ACampo1975}; the preceding chart calculation also shows that its proof is unaffected by branch permutations.

The projection of that boundary fibre to $S$ has, over a point with $r$ local branches, fibre
\[
 \left\{(u_1,\ldots,u_r)\in(S^1)^r:
             \prod_j u_j^{a_j}=e^{\mathrm i\theta}\right\}.
 \tag{C.NC.3}
\]
This is the disjoint union of $\gcd(a_1,\ldots,a_r)$ copies of $(S^1)^{r-1}$: an integral change of basis taking the row $(a_1,\ldots,a_r)$ to $(\gcd(a_j),0,\ldots,0)$ gives the asserted homeomorphism. It has Euler number zero for $r\ge2$ and $a_1$ for $r=1$.

Choose a finite triangulation of the compact analytic strata of $S$, refined so that the normal ray bundles and branch coverings trivialize on every open cell. Over each cell (C.NC.3) is a locally trivial bundle. Additivity over the cells and multiplicativity of compactly supported Euler characteristic give zero for every multiple-branch stratum. On the one-branch stratum of $S_i$ the fibre has $m_i$ points. Normalization identifies that stratum with $T_i\setminus D_i$, giving (C.NC.1). This also proves directly which local contributions enter the global formula.
\end{proof}


% C.2.NC.2 from part_C2_NC.tex

